\chapter{Ensemble Microcanônico}

    O chamado \textit{Ensemble Microcanônico} é definido considerando um valor \textit{fixo} da energia interna de um sistema, $U$, bem como do número de partículas, $N$, e do volume $V$. Procura-se, então, a multiplicidade do macroestado correspondente, $\Omega$ -- isto é, o número de microestados que correspondem à mesma configuração macroscópica. A partir da multiplicidade $\Omega$, podemos obter todas as informações da termodinâmica do sistema a partir da entropia de Boltzmann.

    \begin{statementbox}[label=def:entropia_de_boltzmann]{Entropia de Boltzmann}
        Dado um sistema no macroestado definido pelos valores $(U,V,N)$ e com multiplicidade $\Omega$, a entropia de Boltzmann é a função de estado $S$ definida por:
        \begin{equation}
            S(U,V,N\big)=k_B\ln\Omega,
        \end{equation}
        onde $k_B$ é a constante de Boltzmann, cujo valor exato é $k_B=1,380 649\cdot10^{-23}J/K$.
    \end{statementbox}

     Neste capítulo, vamos resolver três problemas clássicos com a teoria do Ensemble Microcanônico, a saber: o Paramagneto de dois estados, o Sólido de Einstein, e o Gás Ideal Monoatômico.

     \begin{statementbox}{Paramagneto de dois estados}
        O \textbf{Paramagneto de dois estados} consiste em $N$ spins 1/2 que não interagem entre si e sujeitos a um campo magnético uniforme $\mathrm{H}$ aplicado numa dada direção. Cada spin pode ter alinhamento paralelo ao campo, com energia $-\upmu_m\mathrm{H}$, ou antiparalelo ao campo, com energia $+\upmu_m\mathrm{H}$, onde $\upmu_m$ é o momento de dipolo magnético de um spin.
    \end{statementbox}

    \begin{statementbox}{Sólido de Einstein}
        O \textbf{Sólido de Einstein} consiste em $N$ osciladores quânticos \textit{localizados} que não interagem entre si. Cada um pode ter energia $E=n\varepsilon$, com $n=0,1,2,\dots$ e $\varepsilon$ o quanta de energia.
    \end{statementbox}

    \begin{statementbox}{Gás Ideal Monoatômico}
        O \textbf{Gás Ideal Monoatômico} consistem em $N$ partículas não interagentes e indistinguíveis que ocupam um volume $V$ no espaço.
    \end{statementbox}

\section{Encontrando a multiplicidade}

    A principal tarefa no ensemble microcanônico é determinar a multiplicidade $\Omega$ do macrosestado correspondente. Caso o sistema tenha níveis discretos de energia, esta tarefa se resume a responder à pergunta \textit{de quantas formas possíveis podemos distribuir a energia $U$ entre as $N$ partículas?} Em outras palavras, é uma tarefa de análise combinatória.
    
    Sistemas com espectro contínuo de energia representam uma dificuldade maior e são tratados somente na \cref{sec:espec_continuo_de_energia}. Vamos começar com um exemplo simples e seguir incrementando.

    \begin{solved}{Multiplicidade de três partículas de spin 1/2}{multiplicidade_paramagneto_tres_particulas}
        Considere o sistema do paramagnetno de dois estados com energia total fixa $U=-\upmu_m\mathrm{H}$ e formado por $N=3$ spins. Encontre a multiplicidade deste sistema.
        \tcblower
        \noindent\textbf{Resolução:}\\
        \indent Lembrando que cada spin pode estar em dois estados: `para cima' (paralelo ao campo), $\uparrow$, com energia $\varepsilon_{\uparrow}=-\upmu_m\mathrm{H}$, ou `para baixo' (antiparalelo ao campo), $\downarrow$, com energia $\varepsilon_{\downarrow}=\upmu_m\mathrm{H}$.

        Como a energia total é $-\upmu_m\mathrm{H}$, sabemos devemos ter dois spins $\uparrow$ e um spin $\downarrow$, de modo que \[2\varepsilon_{\uparrow}+1\varepsilon_{\downarrow}=-2\upmu_m\mathrm{H}+\upmu_m\mathrm{H}=-\upmu_m\mathrm{H}.\] Resta determinar \textit{de quantas formas} podemos ter dois spins $\uparrow$ e um $\downarrow$. 
        
        Podemos pensar que, inicialmente, todos os três estão no estado $\downarrow$, então escolhemos um spin para colocá-lo no estado $\uparrow$, e depois repetimos com um segundo. Para o primeiro spin, temos 3 opções de escolha, e, para o segundo, temos 2 opções, totalizando $3\cdot2$; contudo, precisamos descontar as permutações possíveis dessa escolha, uma vez que não podemos diferenciar os spins. Concluímos que a multiplicidade é \[\Omega = \frac{\overbrace{3\cdot2}^{\text{Número de escolhas possíveis}}}{\underbrace{2}_{\text{Número de permutações da escolha}}}=3.\]

        Podemos reescrever $\Omega$ numa forma mais interessante. Considerando $N_{\uparrow}=2$ o número de spins no estado $\uparrow$ e $N_{\downarrow}$ o número de spins no estado $\downarrow$, temos: \[\Omega = \frac{3\cdot2}{2}=\frac{3\cdot2\cdot1}{(2\cdot1)\cdot1}=\frac{3!}{2!1!}=\frac{N!}{N_{\uparrow}!N_{\downarrow}!}.\]
    \end{solved}

    \begin{solved}{Multiplicidade do Paramagneto de dois estados}{multiplicidade_paramagneto_dois_estados}
    Considere o sistema do paramagnetno de dois estados com energia total fixa $U$ e formado por $N$ spins. Encontre a multiplicidade deste sistema.
    \tcblower
    \noindent\textbf{Resolução:}\\
    \indent Vamos considerar que o número de spins paralelos ao campo é $N_{\uparrow}$, enquanto o número de spins antiparalelos é $N_{\downarrow}$. Claramente, nós temos \[N=N_{\downarrow}+N_{\uparrow},\] bem como \[U=\upmu_m\mathrm{H}\big(N_{\downarrow}-N_{\uparrow}\big).\] Assim, podemos entender o problema como contar todas as formas de colocarmos $N_{\downarrow}$ dos $N$ spins na posição $\downarrow$, e $N_{\uparrow}$ na posição $\uparrow$.
        
    Para colocar o primeiro spin na posição $\uparrow$, podemos escolher qualquer um dentre as $N$ opções de spin; para o segundo, são $N-1$ opções, e assim por diante, de modo que escolhemos um total de \[\frac{N!}{(N-N_{\uparrow})!}=\frac{N!}{N_{\downarrow}!}\] spins na posição $\uparrow$ (cf. o Ex. Resolvido \ref{solved:multiplicidade_paramagneto_tres_particulas}).
    
    Mas temos que descontar deste número todas as permutações possíveis entre os $N_{\uparrow}$ spins escolhidos, porque não conseguimos distinguir um do outro, e chegamos à nossa multiplicidade: \[\Omega(N_{\uparrow},N)=\frac{N!}{N_{\downarrow}!N_{\uparrow}!}\equiv\binom{N}{N_{\uparrow}}.\]
    Lembrando que $U=\upmu_m\mathrm{H}\big(N_{\downarrow}-N_{\uparrow}\big)$ e $N=N_{\downarrow}+N_{\uparrow}$, podemos escrever $\Omega$ em termos da energia $U$:
        \[\Omega(U,N)=\frac{N!}{\left(\frac{N}{2}+\frac{U}{2\upmu_m\mathrm{H}}\right)!\left(\frac{N}{2}-\frac{U}{2\upmu_m\mathrm{H}}\right)!}.\]
    \end{solved}

    \begin{solved}{Multiplicidade do sólido de Einstein}{}
        Considere o sólido de Einstein com energia total fixa $U$ e formado por $N$ spins. Encontre a multiplicidade deste sistema.
        \tcblower
        \noindent\textbf{Resolução:}\\
        \indent Cada oscilador $j$ tem energia $E_j=n_j\varepsilon$, com $n_j$ natural. A energia total é, então, \[U=E_1+E_2+\cdots+E_N=n_1\varepsilon+n_2\varepsilon+\cdots+n_N\varepsilon=M\varepsilon,\] onde $M$ representa o múltiplo total de energia $\varepsilon$ que devemos distribuir nos $N$ osciladores.

        Para o primeiro oscilador, temos um total de $M$ múltiplos de energia; para o segundo, $M-n_1$, e assim por diante, de modo que temos $M!$ formas de distribuir a energia. Mas como também temos que fixar as posições dos osciladores no espaço, nosso problema se torna a distribuição de $M+N-1$ elementos: $M$ múltiplos de energia e $N-1$ osciladores no espaço, em que o $-1$ aparece porque localização do último oscilador é determinada pela colocação dos $N-1$ anteriores (ele fica no lugar ``que sobra'').

        Levando em conta que tanto os quanta de energia quanto os osciladores são indistinguíveis, concluímos que \[\Omega(U,N)=\frac{(M+N-1)!}{M!(N-1)!},\] que podemos escrever em termos da energia $U=M\varepsilon$: \[\Omega(U,N)=\frac{\left(\frac{U}{\varepsilon}+N-1\right)!}{\left(\frac{U}{\varepsilon}\right)!(N-1)!}.\]
    \end{solved}

    Para tratar o gás ideal, precisamos tomar alguns cuidados. Como o espectro de energia não é discreto, não podemos mais ``contar estados''; por isso, deixamos o tema para a \cref{sec:espec_continuo_de_energia}. Contudo, no exercício abaixo abordamos o problema clássico do gás de rede.

    \begin{solved}{Gás de rede}{}
        Considere um gás de rede constituído por $N$ partículas indistinguíveis distribuídas em $V$ células (com $N\leq V$). Suponha que cada célula possa estar vazia ou ocupada por uma única partícula. Encontre a multiplicidade deste sistema.
        \tcblower
        \noindent\textbf{Resolução:}\\
        \indent O problema consiste em distribuir as $N$ partículas nas $V$ células do sistema. Para isso, a primeira pode `escolher' dentre um total de $V$ células; a segunda, $V-1$ (exclui-se a célula ocupada pela primeira); e assim por diante. Em outras palavras, são
        \[ V\times(V-1)\times(V-2)\times\cdots\times(V-N+1) = \frac{V!}{(V-N)!}\]
        maneiras de ocupar as células. Como as partículas são indistinguíveis, para obtermos a multiplicidade correta basta dividir o nosso resultado por $N!$:
        \begin{equation}\label{eq:multiplicidade_gas_de_rede}
            \Omega(N,V) = \frac{V!}{(V-N)!N!}\,.
        \end{equation}
    \end{solved}

    \begin{statementbox}{Fatoração da multiplicidade}
        Quando um sistema $\mathcal{S}$ pode ser dividido em dois (ou mais) subsistemas independentes, $\mathcal{S}=\mathcal{S}_a\oplus\mathcal{S}_b$, a multiplicidade pode ser fatorada por:
        \begin{equation}
            \Omega_{\mathcal{S}} = \Omega_{a} \times \Omega_{b}\,,
        \end{equation}
        uma vez que, para cada estado do subsistema $\mathcal{S}_a$, existem $\Omega_{b}$ estados disponíveis ao subsistema $\mathcal{S}_b$ (por serem independentes).
    \end{statementbox}

    \begin{solved}{Um modelo de adsorção}{}
        Considere um gás de rede cujas $N$ partículas podem ser adsorvidas em uma superfície. O recipiente disponível para o gás contém $V$ sítios, e a superfície de adsorção, $M$ sítios. As partículas não interagem entre si, mas apenas com a superfície. Assim, uma partícula solta no gás possui energia $U=0$, enquanto uma partícula adsorvida possui energia $U=-\varepsilon$ ($\varepsilon>0$). O sistema está isolado com energia $U=-n\varepsilon$, em que $n\leq N$. Encontre a multiplicidade do sistema.
        \tcblower
        \noindent\textbf{Resolução:}\\
        \indent Uma vez que as partículas do gás não interagem entre si, podemos tratar este problema como dois subsistemas independentes: as partículas do gás no "recipiente", e as partículas adsorvidas. Isso nos permite escrever a multiplicidade como:
        \[ \Omega = \Omega_{\text{livres}}\times\Omega_{\text{adsorvidas}}\,. \]

        Como a energia total é $U=-n\varepsilon$, temos $n$ partículas adsorvidas e, consequentemente, $N-n$ partículas livres. Assim, basta usar a \cref{eq:multiplicidade_gas_de_rede} para o primeiro termo:
        \[ \Omega_{\text{livres}} = \frac{V!}{(V-N-n)!(N-n)!}\,.\]

        Para o segundo, consideramos o problema é análogo: são $M$ sítios de adosorção ao longo dos quais podemos distribuir as $n$ partículas (indistinguíveis) adsorvidas:
        \[ \Omega_{\text{adsorvidas}} = \frac{M!}{(M-n)!n!}\,.\]

        Juntando os dois termos, temos:
        \begin{equation}
            \Omega = \frac{V!}{(V-N-n)!(N-n)!}\times\frac{M!}{(M-n)!n!}\,,
        \end{equation}
        ou, voltando para a variável da energia ($n=-U/\varepsilon$):
        \begin{equation}
            \Omega = \frac{V!}{(V-N+U/\varepsilon)!(N+U/\varepsilon)!}\times\frac{M!}{(M+U/\varepsilon)!(-U/\varepsilon)!}\,.
        \end{equation}

        Com este problema esperamos demonstrar de maneira geral a fundamental propriedade de \textit{fatoração} da multiplicidade.
    \end{solved}

\section{Termodinâmica em ação}

    A partir da multiplicidade, podemos extrair todas as propriedades termodinâmicas do sistema com a entropia de Boltzmann. A Figura \ref{fig:maquina_da_mecanica_estatistica} ilustra o processo: basta \textit{girar a manivela} da mecânica estatística; isto é, transformamos as grandezas estatísticas em potenciais da termodinâmica clássica.

    \begin{figure}[htb]
        \centering
        \resizebox{0.65\linewidth}{!}{\input{figures/stat_mech_machine.tex}}
        \caption{A máquina da mecânica estatística: com as funções de partição num dado ensemble ($\Omega$, $\mathcal{Z}_C$, $\mathcal{Z}_G$), basta ``girarmos'' a manivela do maquinário para extrairmos toda a termodinâmica a partir das informações microscópicas do sistema. Adaptado de blundell (ad. ref).}
        \label{fig:maquina_da_mecanica_estatistica}
    \end{figure}

    \begin{solved}{Entropia do paramagneto de dois estados}{}
        A partir da multiplicidade do paramagneto de dois estados, \[\Omega(U,N)=\frac{N!}{\left(\frac{N}{2}+\frac{U}{2\upmu_m\mathrm{H}}\right)!\left(\frac{N}{2}-\frac{U}{2\upmu_m\mathrm{H}}\right)!},\] determine a entropia de Boltzmann do sistema no limite termodinâmico.
        \tcblower
        \noindent\textbf{Resolução:}\\
        \indent Temos \[S(U,N)=k_B\ln(\Omega)=k_B\left[\ln(N!)+\ln\left(\frac{N}{2}+\frac{U}{2\upmu_m\mathrm{H}}\right)!+\ln\left(\frac{N}{2}-\frac{U}{2\upmu_m\mathrm{H}}\right)!\right].\] No limite termodinâmico, podemos usar a fórmula de Stirling (cf. \cref{eq_appx:approximacao_de_stirling}), com $\ppu=U/N$ fixo:
    \end{solved}

    \begin{solved}{Entropia do sólido de Einstein}{}
        A partir da multiplicidade do sólido de Einstein, \[\Omega(U,N)=\frac{\left(\frac{U}{\varepsilon}+N-1\right)!}{\left(\frac{U}{\varepsilon}\right)!(N-1)!},\] determine a entropia de Boltzmann do sistema no limite termodinâmico.
        \tcblower
        \noindent\textbf{Resolução:}\\
        \indent Temos \[S(U,N)=k_B\ln(\Omega)\]
    \end{solved}

    \begin{solved}{Entropia do gás ideal monoatômico}{}
        A partir da multiplicidade do gás ideal, \[\Omega(U,N)=\frac{3N}{2(3N/2)!}\left(\frac{2\pi m}{h^2}\right)^{\frac{3N}{2}}U^{\frac{3N}{2}-1}\delta U,\] determine a entropia de Boltzmann do sistema no limite termodinâmico.
        \tcblower
        \noindent\textbf{Resolução:}\\
        \indent Temos \[S(U,N)=k_B\ln(\Omega)\]
    \end{solved}

\starredsection{Espectro contínuo de energia\label{sec:espec_continuo_de_energia}}

    Agora voltamos nossa atenção a sistemas com espectro de energia contínuo, isto é, que não pode ser determinado com o uso de múltiplos inteiros de um valor. Nesse caso, a multiplicidade \textbf{não é} determinada com o uso de análise combinatória, mas com uma integral sobre o espaço de fase do sistema. Vamos verificar para dois casos simples e, depois, generalizar o método num terceiro caso mais robusto.

    \begin{statementbox}{Multiplicidade com espectro contínuo de energia}
        Dado um sistema com $N$ graus de liberdade e energia $U$ fixa, integramos sobre todas as posições e momentos acessíveis ao sistema:
        \[ \Omega(U,N)=\frac{1}{h^{N}}\int\int_{\text{Energia }U} \diff \X_1\diff \X_2\cdots\diff \X_N\diff \p_1\diff \p_2\cdots\diff \p_N,\] onde $h$ é a constante de Planck.
    \end{statementbox}

    \begin{solved}{Multiplicidade do oscilador harmônico unidimensional}{}
        Considere um único oscilador harmônico unidimensional de massa $m$ e frequência $\omega$ tem energia fixa, dada por \[U=\frac{\p^2}{2m}+\frac{1}{2}m\omega^2\X^2.\] Determine a multiplicidade deste sistema.
        \tcblower
        \noindent\textbf{Resolução:}\\
        \indent Começamos notando que \[U=\frac{\p^2}{2m}+\frac{1}{2}m\omega^2\X^2\] é uma \textit{equação de elipse} para as variáveis $\p$ e $\X$. Desse modo, a multiplicidade $\Omega$ será dada por uma coroa elíptica, por um pequeno desvio $\delta U$ da energia. 
        
        A Figura \ref{fig:oscilador_harmonico_espaco_de_fase} vai nos guiar neste cálculo. A área da elipse maior, com energia $U+\delta U$, é \[\pi\sqrt{\frac{2(U+\delta U)}{m\omega^2}}
\sqrt{2m(U+\delta U)} = \pi\frac{2(U+\delta U)}{\omega},\] e a área da elipse menor, é \[\pi\frac{2U}{\omega},\] de modo que \[\Omega = \frac{1}{h}\left(\pi\frac{2(U+\delta U)}{\omega}-\pi\frac{2U}{\omega} \right)=\frac{2\pi}{h\omega}\delta U.\] Vemos que, neste caso muito simples, o volume do espaço de fase acessível ao sistema é uma função independente da energia!

        \begin{center}
            % 1. Coloca a âncora do link no topo e define o tipo
            \captionsetup{type=figure} 
            
            % 2. Insere o seu gráfico TikZ
            \begin{tikzpicture}[scale=1]

    % 1. Preenchimento (O "Volume" no espaço de fase / Multiplicidade)
    % Preenchemos primeiro para que as linhas fiquem por cima
    \fill[ThemeColor!10!white, even odd rule]
        %> Elipse energia U
        (0,0) ellipse (3.5cm and 2.2cm)
        (0,0) ellipse (3.1cm and 1.85cm);

    % 2. A curva da elipse (Trajetória de energia constante U)
    \draw[very thick, ThemeColor] (0,0) ellipse (3.5cm and 2.2cm) (0,0) ellipse (3.1cm and 1.85cm);

    % 3. Eixos do Espaço de Fase
    \draw[->, thick] (-4.5,0) -- (4.5,0) node[right] {$\X$};
    \draw[->, thick] (0,-3.2) -- (0,3.2) node[above] {$\p$};

    % 4. Marcações (Limites de x e p baseados na equação)
    % Limites no eixo x (quando p = 0)
    \draw (3.5,0.1) -- (3.5,-0.1) node[below right] {$\sqrt{\frac{2U}{m\omega^2}}$};
    \draw (-3.5,0.1) -- (-3.5,-0.1) node[below left] {$-\sqrt{\frac{2U}{m\omega^2}}$};
    
    % Limites no eixo p (quando x = 0)
    \draw (0.1,2.2) -- (-0.1,2.2) node[above left] {$\sqrt{2mU}$};
    \draw (0.1,-2.2) -- (-0.1,-2.2) node[below left] {$-\sqrt{2mU}$};

\end{tikzpicture}
            
            % 3. Agora você volta a usar o \caption normal (sem o 'of')
            \caption{Espaço de fase de um oscilador harmônico unidimensional.}
            \label{fig:oscilador_harmonico_espaco_de_fase}
        \end{center}
    \end{solved}

    \begin{solved}{Multiplicidade de uma partícula de gás ideal}{}
        Considere uma partícula com massa $m$ e energia \[U=\frac{\vec{\p}^2}{2m},\] contida num volume $V$. Determine a multiplicidade deste sistema.
        \tcblower
        \noindent\textbf{Resolução:}\\
        \indent Neste caso, a posição $\vec{\R}$ da partícula é \textit{independente} do seu momento, $\vec{\p}$, de modo que podemos escrever a multiplicidade como \[\Omega = \frac{1}{h^3}\int_V\diff^3\R\int_{U=\vec{\p}^2/2m}\diff^3\p=\frac{V}{h^3}\int_{U=\vec{\p}^2/2m}\diff^3\p.\]
        Assim, vamos obter a multiplicidade pelo volume de uma coroa esférica por um pequeno desvio $\delta U$ da energia.

        Como $\p^2=2mU$, o raio da esfera maior é $\sqrt{2m(U+\delta U)}$, e o da menor, $\sqrt{2mU}$. Assim: \[\frac{4\pi}{3}\big(2m(U+\delta U)\big)^{3/2}-\frac{4\pi}{3}\big(2mU\big)^{3/2}=\frac{8\pi}{3}m^{3/2}U^{3/2}\left[\left(1+\frac{\delta U}{U}\right)^{3/2}-1\right],\] usando a série binomial (cf. Eq. \eqref{eq_appx:serie_binomial}), \[\frac{8\pi}{3}m^{3/2}U^{3/2}\left[\left(1+\frac{\delta U}{U}\right)^{3/2}-1\right]\approx\frac{8\pi}{3}m^{3/2}U^{3/2}\left[1+\frac{3}{2}\frac{\delta U}{U}-1\right]=4\pi m^{3/2}U^{1/2}\delta U.\] Finalmente, encontramos \[\Omega = \frac{4\pi m^{3/2}V}{h^3}U^{1/2}\delta U.\]
    \end{solved}

    \begin{solved}{Multiplicidade do Gás Ideal Monoatômico}{}
        Considere um gás ideal composto por $N$ partículas idênticas de massa $m$, cada uma restringida a se mover num volume $V$ fixo, com energia $U$ dada por \[U=\frac{\vec{\p}^2_1}{2m}+\frac{\vec{\p}^2_2}{2m}+\cdots+\frac{\vec{\p}^2_N}{2m}=\sum_{j=1}^N\frac{\vec{\p}^2_j}{2m}.\] Determine a multiplicidade desse sistema.
        \tcblower
        \noindent\textbf{Resolução:}\\
        \indent Procedemos como no exercício anterior: \[\Omega(U,N)=\frac{1}{h^{3N}}\idotsint_V\diff^3\R_1\cdots\diff^3\R_N\idotsint_{U=\sum\vec{\p}^2_j/2m}\diff^3\p_1\dots\diff^3\p_N.\] A integral sobre a posição de cada uma das partículas leva a $V$, resultando, para as $N$ integrais, em $V^N$. Já a integral sobre os $N$ momentos, utilizaremos o método anterior: corresponde ao volume de uma coroa esférica por um pequeno desvio $\delta U$ da energia.

        Usando a Eq. \eqref{eq_appx:volume_hiperesfera} para o volume de uma hiperesfera de raio $R$, temos: \[V_{3N}(\sqrt{2m(U+\delta U)})=\frac{(2\pi m)^{3N/2}}{\big(3N/2\big)!}(U+\delta U)^{3N/2}\underbrace{\approx}_{\text{Série binomial}}\frac{(2\pi mU)^{3N/2}}{\big(3N/2\big)!}\left(1+\frac{3N}{2}\frac{\delta U}{U}\right),\]
        de modo que a integral sobre os momentos fica: \[\idotsint_{U=\sum\vec{\p}^2_j/2m}\diff^3\p_1\dots\diff^3\p_N = \frac{3N}{2(3N/2)!}\big(2\pi m\big)^{3N/2}U^{3N/2-1}\delta U,\] resultando em:
        \[\Omega(U,N)=\frac{3N}{2(3N/2)!}\left(\frac{2\pi m}{h^2}\right)^{\frac{3N}{2}}U^{\frac{3N}{2}-1}\delta U.\] Contudo, entra aqui um problema: precisamos adicionar um fator \textit{ad hoc} a este termo para que o limite termodinâmicos corresponda à correta representação do gás ideal. Comentamos sobre isto no Ex. Resolvido ADICIONAR EXERCÍCIO RESOLVIDO.
    \end{solved}

\newpage
\section*{Exercícios Propostos}
\begin{questionsbox}[Sistemas com energia discreta]
    \begin{question} %> Callen
    Considere um sistema de $N$ partículas livres que podem assumir apenas seus estados de energia $0$ e $\varepsilon$ (com $\varepsilon>0$). Sejam $n_0$ e $n_1$ os números de ocupação dos níveis de energia $0$ e $\varepsilon$, respectivamente. A energia total do sistema é $U$.
    \begin{parts}
        \item Determine a entropia deste sistema.
        \item Obtenha a temperatura como uma função de $U$ e mostre que ela pode ser negativa.
        \item O que acontece quando um sistema com temperatura negativa pode trocar calor com um sistema com temperatura positiva?
    \end{parts}
    \end{question}
    \begin{answer}
        \begin{parts}
        \item $\pps(\ppu) = -k\frac{\ppu}{\varepsilon}\ln\big(\varepsilon/\ppu-1\big)+k\ln\big(1-\ppu/\varepsilon\big)$
        \item $1/T=-k/\varepsilon\ln\big(\varepsilon/\ppu-1\big)$
    \end{parts}
    \end{answer}
\end{questionsbox}

\section*{Problemas Propostos}

\begin{problem}[subtitle={Teorema do retorno de Poicaré}]

    
\end{problem} 